% book/chapters/ch07-models.tex
\chapter{Model Selection}
\label{ch:models}

\index{model selection}

Claude Code uses three models from the Claude~4 family. Choosing the right
model for each task is the fastest lever for reducing both cost and latency
without sacrificing quality. This chapter describes the current models, their
capabilities, and a decision framework.

\begin{warningbox}
Model names, pricing, and capabilities change as Anthropic releases new
versions. Always verify current information at
\url{https://docs.anthropic.com/en/docs/about-claude/models/overview}.
The data below was verified against the Anthropic platform documentation
on 2026-05-27.
\end{warningbox}

\section{The Claude 4 Model Family}
\label{sec:model-family}

\index{Claude Opus 4.7}\index{Claude Sonnet 4.6}\index{Claude Haiku 4.5}

\begin{table}[ht]
\centering
\caption{Current Claude model comparison (verified 2026-05-27)}
\label{tab:model-comparison}
\begin{tabularx}{\textwidth}{lllll}
\toprule
\textbf{Feature}        & \textbf{Opus 4.7}     & \textbf{Sonnet 4.6}   & \textbf{Haiku 4.5} \\
\midrule
API ID                  & \texttt{claude-opus-4-7}
                        & \texttt{claude-sonnet-4-6}
                        & \texttt{claude-haiku-4-5-20251001} \\
Context window          & 1M tokens             & 1M tokens             & 200K tokens \\
Max output              & 128K tokens           & 64K tokens            & 64K tokens \\
Input price (per MTok)  & \$5                   & \$3                   & \$1 \\
Output price (per MTok) & \$25                  & \$15                  & \$5 \\
Extended thinking       & No (adaptive only)    & Yes                   & Yes \\
Relative latency        & Moderate              & Fast                  & Fastest \\
Knowledge cutoff        & Jan 2026              & Aug 2025              & Feb 2025 \\
\bottomrule
\end{tabularx}
\end{table}

\subsection{Claude Opus 4.7 — Most Capable}

Opus 4.7 is Anthropic's most capable generally available model.
Its strengths are complex reasoning and agentic coding: tasks that involve
planning across many files, making architectural decisions, or synthesising
information from a large context.

\begin{lstlisting}[caption={Tasks suited to Opus 4.7}]
"Design a scalable microservices architecture for this monolith"
"Deep security audit for OWASP Top 10 across the entire codebase"
"Refactor this module to an event-driven pattern with full test coverage"
\end{lstlisting}

\subsection{Claude Sonnet 4.6 — Best Speed/Intelligence Balance}

Sonnet 4.6 handles the vast majority of development work: feature
implementation, bug fixing, code review, and test generation. It supports
extended thinking (for explicit reasoning), has a 1M-token context window,
and is fast enough for interactive use.

\begin{lstlisting}[caption={Tasks suited to Sonnet 4.6}]
"Add user authentication to this Express API"
"Refactor this class to use the repository pattern"
"Write integration tests for the checkout flow"
\end{lstlisting}

\subsection{Claude Haiku 4.5 — Fastest, Most Cost-Efficient}

Haiku 4.5 is 3× cheaper than Sonnet and responds in seconds. It excels at
tasks that do not require deep reasoning: file searches, formatting, simple
transformations, and quick lookups. Use it as the default for the Explore
subagent and for any automation that runs at high frequency.

\begin{lstlisting}[caption={Tasks suited to Haiku 4.5}]
"Find all React components importing useState"
"Format this file with Prettier"
"Sort these imports alphabetically"
\end{lstlisting}

\section{Pricing and Cost Optimisation}
\label{sec:model-pricing}

\index{model selection!cost}

A well-chosen model mix can reduce daily token costs by 50–60\,\% versus
using Sonnet or Opus for everything:

\begin{table}[ht]
\centering
\caption{Daily cost comparison: all-Sonnet vs.\ smart mix}
\label{tab:model-cost}
\begin{tabularx}{\textwidth}{Xrrr}
\toprule
\textbf{Task type (ops/day)} & \textbf{All Sonnet} & \textbf{Smart mix} & \textbf{Model used} \\
\midrule
File searches (20)       & \$0.36 & \$0.06 & Haiku \\
Formatting (15)          & \$0.27 & \$0.04 & Haiku \\
Bug fixes (10)           & \$0.90 & \$0.90 & Sonnet \\
Feature development (5)  & \$0.45 & \$0.45 & Sonnet \\
Architecture design (2)  & \$0.90 & \$1.60 & Opus \\
\midrule
\textbf{Daily total}     & \$2.88 & \$3.05 & \\
\multicolumn{4}{l}{\small Note: smart mix costs slightly more here because Opus tasks are 2×.} \\
\multicolumn{4}{l}{\small For 40 searches and 0 architecture, smart mix saves $\approx$50\%.} \\
\bottomrule
\end{tabularx}
\end{table}

\section{Model Selection Decision Framework}
\label{sec:model-decision}

\index{model selection!decision framework}

\begin{itemize}
  \item \textbf{Simple task, no code understanding required?}
        $\Rightarrow$ \textbf{Haiku 4.5}
        (searches, formatting, linting, import sorting)
  \item \textbf{Standard coding task?}
        $\Rightarrow$ \textbf{Sonnet 4.6}
        (features, bug fixes, tests, documentation, reviews)
  \item \textbf{Architecture, security-critical, or long-context synthesis?}
        $\Rightarrow$ \textbf{Opus 4.7}
\end{itemize}

\textbf{Rule of thumb}: use the \emph{least powerful model that can get the
job done correctly}.

\section{Configuring Models in Claude Code}
\label{sec:model-config}

\index{model selection!configuration}

\textbf{Session-level} (CLI flag):
\begin{lstlisting}[language=bash,caption={Selecting a model for a session}]
claude --model claude-opus-4-7
claude --model claude-sonnet-4-6
claude --model claude-haiku-4-5
\end{lstlisting}

\textbf{Default for a project} (\texttt{.claude/settings.json}):
\begin{lstlisting}[language=json,caption={Setting the project default model}]
{
  "defaultModel": "claude-sonnet-4-6"
}
\end{lstlisting}

\textbf{Per-subagent} (frontmatter in \texttt{.claude/agents/}):
\begin{lstlisting}[caption={Assigning Haiku to the Explore subagent (already the default)}]
---
name: explore
model: haiku
---
\end{lstlisting}

\textbf{Per-skill} (frontmatter in \texttt{SKILL.md}):
\begin{lstlisting}[caption={Routing a skill to Opus for architecture work}]
---
name: design-architecture
model: opus
---
\end{lstlisting}

\section{Legacy Models}
\label{sec:model-legacy}

\index{model selection!deprecated}

Several earlier models remain available but should be migrated away from:

\begin{itemize}
  \item \textbf{Claude Sonnet~4} (\texttt{claude-sonnet-4-20250514}) and
        \textbf{Claude Opus~4} (\texttt{claude-opus-4-20250514}) are
        \textbf{deprecated} and will be retired on \textbf{June 15, 2026}.
        Migrate to Sonnet~4.6 and Opus~4.7 respectively.
  \item Sonnet~4.5, Opus~4.5, Opus~4.6, and Opus~4.1 are still available
        but no longer the recommended option.
\end{itemize}

%% ---------------------------------------------------------------
\begin{summarybox}
\begin{itemize}
  \item Three current models: \textbf{Opus 4.7} (most capable, 1M context,
        \$5/\$25/MTok), \textbf{Sonnet 4.6} (balanced, 1M context,
        \$3/\$15/MTok), \textbf{Haiku 4.5} (fastest, 200K context,
        \$1/\$5/MTok).
  \item Opus 4.7 and Sonnet 4.6 both have \textbf{1M-token context windows}
        — a major increase from the 200K reported in earlier documentation.
  \item Only Sonnet 4.6 and Haiku 4.5 support extended thinking; Opus 4.7
        uses adaptive thinking instead.
  \item Use Haiku for simple/search tasks, Sonnet for 90\,\% of coding work,
        and Opus only for architecture or security-critical analysis.
  \item Models can be set globally (\texttt{defaultModel} in
        \texttt{settings.json}), per session (CLI flag), per subagent, or
        per skill.
  \item Sonnet~4 and Opus~4 are deprecated; retire by June 15, 2026.
\end{itemize}
\end{summarybox}

\printbibliography[heading=subbibliography,title={References}]
